\chapter{Proof of the Fundamental Theorem}\label{ch:assembly}

This chapter proves the Fundamental Theorem of Matrix Product States.
We state the Fundamental Theorem in two versions: one for block-injective
tensors in canonical form with basis of normal tensors (BNT) separation
(CF-BNT), and one for the normal-canonical variant.
We then treat the equal-MPV case, including the common-block-structure
specialization that deduces gauge equivalence.
The results combine the canonical form reduction
(Chapter~\ref{ch:canonical}),
block separation (Chapter~\ref{ch:block_perm}),
spectral gap (Chapter~\ref{ch:spectral}),
and the BNT permutation arguments of Chapter~\ref{ch:bnt}.

The theorems in Sections~\ref{sec:assembly_weak}--\ref{sec:assembly_equal}
follow the approach of \cite{Cirac2017MPDO_AnnPhys}:
they first block the tensor to remove non-trivial periodicity, then
apply the aperiodic fundamental theorem.
Section~\ref{sec:periodic_ft} outlines the stronger
``irreducible form'' framework of \cite{DeLasCuevas2017Irreducible},
which gives a fundamental theorem valid for periodic tensors directly,
with an extra $Z$-gauge degree of freedom.

\section{Weak Fundamental Theorem}%
\label{sec:assembly_weak}

The following theorem combines the full upstream reduction
(zero-block separation, TP gauge, blocking, primitivity)
with the downstream Fundamental Theorem for normal canonical form data.
It is conditional on the extra hypotheses (tensor irreducibility and
distinct weight norms) that the automatic reduction does not yet produce.

\begin{theorem}[Weak Fundamental Theorem (conditional)]%
    \label{thm:weak_ft_conditional}
    \lean{MPSTensor.weakFundamentalTheorem_conditional}
    \leanok
    \uses{def:is_normal_canonical_form, def:gauge_phase_equiv}
    Let $(\mu_j^A, A_j)_{j=1}^{g_A}$ and $(\mu_k^B, B_k)_{k=1}^{g_B}$
    be two families in normal canonical form,
    each satisfying the non-degeneracy condition
    that no two distinct blocks in the same family are gauge-phase
    equivalent.
    Let $A_\mathrm{tot}$ and $B_\mathrm{tot}$ be tensors with
    MPV decompositions
    \[
        V^{(N)}(A_\mathrm{tot})_\sigma = \sum_j a_{N,j}\, V^{(N)}(A_j)_\sigma,
        \quad
        V^{(N)}(B_\mathrm{tot})_\sigma = \sum_k b_{N,k}\, V^{(N)}(B_k)_\sigma,
    \]
    with $a_{N,j} \to a_j \neq 0$ and $b_{N,k} \to b_k \neq 0$.
    Suppose moreover that
    $V^{(N)}(A_\mathrm{tot})_\sigma = c_N\, V^{(N)}(B_\mathrm{tot})_\sigma$
    with $c_N \to c \neq 0$.
    Then $g_A = g_B$ and there is a permutation $\pi$ such that
    $D_j^A = D_{\pi(j)}^B$ and $A_j$ is gauge-phase equivalent to
    $B_{\pi(j)}$ for every~$j$.
\end{theorem}

\begin{proof}\leanok
    \uses{thm:ft_proportional_nt}
    This packages the normal-canonical-form Fundamental Theorem
    (Theorem~\ref{thm:ft_proportional_nt}).
    The normal canonical form data provide the BNT separation
    and overlap convergence hypotheses;
    the non-degeneracy condition ensures the block matching
    is well-defined.
\end{proof}

\section{Proportional Fundamental Theorem}%
\label{sec:assembly_proportional}

\begin{theorem}[Fundamental Theorem --- proportional MPV case]\label{thm:ft_proportional}
	\lean{MPSTensor.fundamentalTheorem_proportionalMPV_CFBNT}
	\leanok
	\uses{def:cf_bnt, def:gauge_phase_equiv}
	Let $A_{\mathrm{tot}}$ and $B_{\mathrm{tot}}$ be MPS tensors
	in canonical form with BNT separation, with associated families
	$(\mu_j^A, A_j)_{j=1}^{g_A}$ and $(\mu_k^B, B_k)_{k=1}^{g_B}$.
	Suppose there exist coefficient arrays $a_{N,j}$, $b_{N,k}$ and
	nonzero limits $a_j^\infty$, $b_k^\infty$, $c^\infty$ such that
	\[
		V^{(N)}(A_{\mathrm{tot}})_\sigma
		= \sum_{j=1}^{g_A} a_{N,j} V^{(N)}(A_j)_\sigma,
		\qquad
		V^{(N)}(B_{\mathrm{tot}})_\sigma
		= \sum_{k=1}^{g_B} b_{N,k} V^{(N)}(B_k)_\sigma,
	\]
	with $a_{N,j} \to a_j^\infty \neq 0$ and
	$b_{N,k} \to b_k^\infty \neq 0$, and suppose also that
	\[
		V^{(N)}(A_{\mathrm{tot}})_\sigma
		= c_N V^{(N)}(B_{\mathrm{tot}})_\sigma
	\]
	for a sequence $c_N \to c^\infty \neq 0$.
	Then $g_A = g_B$, and there is a permutation $\pi$ such that
	$D_j^A = D_{\pi(j)}^B$ and $A_j$ is gauge-phase equivalent to
	$B_{\pi(j)}$ for each~$j$.
\end{theorem}

\begin{proof}
	\leanok
	\uses{thm:cross_overlap_zero, thm:bnt_perm_prop}
	Apply Theorem~\ref{thm:bnt_perm_prop} to the two CF-BNT
	decompositions.
	The CF-BNT hypotheses provide injectivity, TP normalization,
	diagonal overlap limits, and off-diagonal decay via
	Theorem~\ref{thm:cross_overlap_zero}.
	With the supplied decomposition identities and convergence
	hypotheses, Theorem~\ref{thm:bnt_perm_prop} yields the equality of
	block counts, the permutation, and the per-block gauge-phase
	equivalence.
\end{proof}

\begin{remark}[Coefficient arrays in the CF-BNT setting]
	In applications the coefficient arrays in
	Theorem~\ref{thm:ft_proportional} come from the block-diagonal
	decomposition itself. If
	\[
		A_{\mathrm{tot}} = \bigoplus_{j=1}^{g_A} \mu_j^A A_j,
		\qquad
		B_{\mathrm{tot}} = \bigoplus_{k=1}^{g_B} \mu_k^B B_k,
	\]
	then Theorem~\ref{thm:mpv_decomposition} gives
	\[
		V^{(N)}(A_{\mathrm{tot}})_\sigma
		= \sum_{j=1}^{g_A} (\mu_j^A)^N V^{(N)}(A_j)_\sigma,
		\qquad
		V^{(N)}(B_{\mathrm{tot}})_\sigma
		= \sum_{k=1}^{g_B} (\mu_k^B)^N V^{(N)}(B_k)_\sigma.
	\]
	Thus one takes $a_{N,j} = (\mu_j^A)^N$ and
	$b_{N,k} = (\mu_k^B)^N$, while the sequence $c_N$ is the global
	proportionality factor from the hypothesis
	$V^{(N)}(A_{\mathrm{tot}}) = c_N V^{(N)}(B_{\mathrm{tot}})$.
	Theorem~\ref{thm:ft_proportional} keeps these arrays explicit because
	Theorem~\ref{thm:bnt_perm_prop} is formulated in that abstract form.
\end{remark}

\begin{theorem}[Fundamental Theorem --- proportional MPV case, NT route]\label{thm:ft_proportional_nt}
		\lean{MPSTensor.fundamentalTheorem_proportionalMPV_normalCFBNT}\leanok
		\uses{def:normal_cf_bnt, def:gauge_phase_equiv}
		Let $A_{\mathrm{tot}}$ and $B_{\mathrm{tot}}$ be MPS tensors
		in normal canonical form with BNT separation, with associated families
		$(\mu_j^A, A_j)_{j=1}^{g_A}$ and $(\mu_k^B, B_k)_{k=1}^{g_B}$.
		Suppose there exist coefficient arrays $a_{N,j}$, $b_{N,k}$ and
		nonzero limits $a_j^\infty$, $b_k^\infty$, $c^\infty$ such that
		\[
			V^{(N)}(A_{\mathrm{tot}})_\sigma
			= \sum_{j=1}^{g_A} a_{N,j} V^{(N)}(A_j)_\sigma,
			\qquad
			V^{(N)}(B_{\mathrm{tot}})_\sigma
			= \sum_{k=1}^{g_B} b_{N,k} V^{(N)}(B_k)_\sigma,
		\]
		with $a_{N,j} \to a_j^\infty \neq 0$ and
		$b_{N,k} \to b_k^\infty \neq 0$, and suppose also that
		\[
			V^{(N)}(A_{\mathrm{tot}})_\sigma
			= c_N V^{(N)}(B_{\mathrm{tot}})_\sigma
		\]
		for a sequence $c_N \to c^\infty \neq 0$.
		Then $g_A = g_B$, and there is a permutation $\pi$ such that
		$D_j^A = D_{\pi(j)}^B$ and $A_j$ is gauge-phase equivalent to
		$B_{\pi(j)}$ for each~$j$.
\end{theorem}

\begin{proof}\leanok
		\uses{thm:ncf_overlap_tendsto_one, thm:nt_modulus_one_gauge, rem:normalcf_bnt_gap}
		Normal canonical form with BNT separation provides irreducibility,
		TP normalization, and primitive self-overlap convergence via
		Theorem~\ref{thm:ncf_overlap_tendsto_one}.
		Theorem~\ref{thm:nt_modulus_one_gauge} supplies the leading-block
		identification step: a nondecaying mixed overlap forces the
		corresponding comparison block to be gauge-phase equivalent.
		The same peeling argument as in the injective CF-BNT case therefore
		yields the block permutation and gauge-phase equivalence.
		Remark~\ref{rem:normalcf_bnt_gap} explains why the argument stops at
		this separated normal-canonical level.
\end{proof}

\section{Equal-MPV Fundamental Theorem}%
\label{sec:assembly_equal}

\begin{theorem}[Fundamental Theorem --- equal MPV case (target statement)]\label{thm:ft_equal}
	\notready
	\uses{def:cf_bnt, def:gauge_equiv}
	Let $A_{\mathrm{tot}}$ and $B_{\mathrm{tot}}$ be MPS tensors
	in canonical form with BNT separation, with associated families
	$(\mu_j^A, A_j)_{j=1}^{g_A}$ and $(\mu_k^B, B_k)_{k=1}^{g_B}$.
	If
	\[
		V^{(N)}(A_{\mathrm{tot}})_\sigma
		= V^{(N)}(B_{\mathrm{tot}})_\sigma
	\]
	for every $N$ and $\sigma$, then the block-diagonal tensors
	$\bigoplus_{j=1}^{g_A} \mu_j^A A_j$ and
	$\bigoplus_{k=1}^{g_B} \mu_k^B B_k$
	are gauge equivalent.
\end{theorem}

\begin{proof}
	\uses{thm:ft_proportional, thm:mpv_decomposition, thm:power_sum_multiset}
	The intended literature route begins by applying
	Theorem~\ref{thm:ft_proportional} with $c_N = 1$, which yields a
	permutation of the matched blocks and blockwise gauge-phase equivalences.
	One then compares the resulting MPV expansions from
	Theorem~\ref{thm:mpv_decomposition}. In the general equal case the
	weights occur with multiplicities, so recovering the individual weights
	requires a power-sum argument such as
	Theorem~\ref{thm:power_sum_multiset}. After the weights are identified,
	the phase factors are absorbed into them and the blockwise conjugacies
	assemble to a global gauge equivalence. This full route is not yet
	formalized here.
\end{proof}

\begin{theorem}[Fundamental Theorem --- equal MPV case with explicit coefficient data]\label{thm:ft_equal_explicit}
	\lean{MPSTensor.fundamentalTheorem_equalMPV_full}
	\leanok
	\uses{def:cf_bnt, def:gauge_equiv}
	Let $A_{\mathrm{tot}}$ and $B_{\mathrm{tot}}$ be MPS tensors
	in canonical form with BNT separation, with associated families
	$(\mu_j^A, A_j)_{j=1}^{g_A}$ and $(\mu_k^B, B_k)_{k=1}^{g_B}$.
	Suppose there exist coefficient arrays $a_{N,j}$, $b_{N,k}$ and
	nonzero limits $a_j^\infty$, $b_k^\infty$ such that
	\[
		V^{(N)}(A_{\mathrm{tot}})_\sigma
		= \sum_{j=1}^{g_A} a_{N,j} V^{(N)}(A_j)_\sigma,
		\qquad
		V^{(N)}(B_{\mathrm{tot}})_\sigma
		= \sum_{k=1}^{g_B} b_{N,k} V^{(N)}(B_k)_\sigma,
	\]
	with $a_{N,j} \to a_j^\infty \neq 0$ and
	$b_{N,k} \to b_k^\infty \neq 0$.
	If
	\[
		V^{(N)}(A_{\mathrm{tot}})_\sigma
		= V^{(N)}(B_{\mathrm{tot}})_\sigma
	\]
	for every $N$ and $\sigma$, then $g_A = g_B$, and there is a
	permutation $\pi$ such that $D_j^A = D_{\pi(j)}^B$ and, after
	reindexing the $B$-blocks by $\pi$, the weighted block-diagonal tensors
	$\bigoplus_{j=1}^{g_A} \mu_j^A A_j$ and
	$\bigoplus_{j=1}^{g_A} \mu_{\pi(j)}^B B_{\pi(j)}$
	are gauge equivalent.
\end{theorem}

\begin{proof}
	\leanok
	\uses{thm:ft_proportional, thm:cf_bnt_is_bnt, def:bnt, thm:mpv_decomposition, def:gauge_phase_equiv}
	Apply Theorem~\ref{thm:ft_proportional} with $c_N = 1$.
	This yields $g_A = g_B$, a permutation $\pi$, equal block dimensions,
	and for each $j$ a gauge-phase equivalence
	\[
		B_{\pi(j)}^i = \zeta_j X_j A_j^i X_j^{-1}.
	\]
	By Theorem~\ref{thm:mpv_decomposition},
	\[
		V^{(N)}(A_{\mathrm{tot}})_\sigma
		= \sum_{j=1}^{g_A} (\mu_j^A)^N V^{(N)}(A_j)_\sigma,
	\]
	and similarly for $B_{\mathrm{tot}}$.
	Reindex the $B$-sum by $\pi$ and use the gauge-phase identity to rewrite
	\[
		V^{(N)}(B_{\pi(j)})_\sigma = \zeta_j^N V^{(N)}(A_j)_\sigma.
	\]
	Then equality of MPVs implies
	\[
		\sum_{j=1}^{g_A}
		\left((\mu_j^A)^N - (\mu_{\pi(j)}^B \zeta_j)^N\right)
		V^{(N)}(A_j)_\sigma = 0.
	\]
	By Theorem~\ref{thm:cf_bnt_is_bnt} and Definition~\ref{def:bnt}, the
	family $\{V^{(N)}(A_j)\}_{j=1}^{g_A}$ is linearly independent for all
	sufficiently large $N$, so
	\[
		(\mu_j^A)^N = (\mu_{\pi(j)}^B \zeta_j)^N
	\]
	for every $j$ and all large $N$.
	Comparing two consecutive values of $N$ and using that the weights are
	nonzero gives
	\[
		\mu_j^A = \mu_{\pi(j)}^B \zeta_j.
	\]
	Thus the phase factors can be absorbed into the weights, and the
	blockwise conjugacies assemble to a gauge equivalence of the reindexed
	weighted block-diagonal tensors.
\end{proof}

\begin{theorem}[Fundamental Theorem --- equal MPV case, common block structure]\label{thm:ft_equal_same_structure}
	\lean{MPSTensor.fundamentalTheorem_equalMPV_CFBNT}
	\leanok
	\uses{def:cf_bnt, def:gauge_equiv}
	Fix a common block structure: the same number of blocks~$g$,
	the same bond dimensions $(D_k)_{k=1}^g$, and the same weights
	$(\mu_k)_{k=1}^g$.
	Let $(\mu_k, A_k)_{k=1}^g$ and $(\mu_k, B_k)_{k=1}^g$ be two families in
	canonical form with BNT separation on this common data.
	If the block-diagonal tensors
	$\bigoplus_k \mu_k A_k$ and $\bigoplus_k \mu_k B_k$
	generate the same MPV family, then each pair $(A_k, B_k)$ is
	gauge equivalent, and hence the assembled block-diagonal tensors
	are gauge equivalent.
\end{theorem}

\begin{proof}
	\leanok
	\uses{thm:ft_canonical}
	Since both families share the canonical form data, the BNT separation
	condition is not needed.
	The canonical-form fundamental theorem
	(Theorem~\ref{thm:ft_canonical}) then applies directly to this
	common $(g, D_k, \mu_k)$ structure, yielding per-block gauge
	equivalence and hence global gauge equivalence.
\end{proof}

\subsection{Route~B: Paper-faithful multiplicity layer}%
\label{sec:route_b}

The paper-faithful multiplicity interface keeps the literature's
per-block-copy structure separate from the merged
\texttt{IsCanonicalFormBNT} predicate.

\begin{definition}[Paper sector data and decomposition]%
\label{def:paper_sector_data}
	\lean{MPSTensor.SectorWeightData, MPSTensor.SectorDecomposition}
	\leanok
	A \emph{paper sector data} over $g$ basis blocks specifies:
	multiplicities $r_j > 0$, sector weights
	$\mu_{j,q} \neq 0$ for $q \in \{1,\dotsc,r_j\}$, and the
	\emph{paper coefficient}
	$\mathrm{coeff}(N,j) := \sum_{q=1}^{r_j} \mu_{j,q}^N$.
	A \emph{paper decomposition} bundles the basis blocks $(A_j)_{j=1}^g$
	with their sector data.
\end{definition}

\begin{theorem}[Paper decomposition formula]%
\label{thm:paper_decomp_formula}
	\lean{MPSTensor.SectorDecomposition.mpv_toTensor_eq_sum_coeff}
	\leanok
	\uses{def:paper_sector_data}
	For a paper decomposition $P$, the MPV of the assembled tensor satisfies
	\[
		V^{(N)}(P.\mathrm{toTensor})_\sigma
		= \sum_{j=1}^{g} \mathrm{coeff}(N,j) \cdot V^{(N)}(A_j)_\sigma.
	\]
\end{theorem}

\begin{theorem}[Eventual coefficient comparison from BNT linear independence]%
\label{thm:coeff_eventually_eq}
	\lean{MPSTensor.SectorWeightData.coeff_eventually_eq_of_sameMPV}
	\leanok
	\uses{def:paper_sector_data, def:bnt}
	If two sector data $S, T$ over the same basis produce equal total MPVs
	and the basis is eventually linearly independent (the BNT property),
	then $S.\mathrm{coeff}(N,j) = T.\mathrm{coeff}(N,j)$ for all
	sufficiently large $N$.
\end{theorem}

\begin{theorem}[Geometric sequence extrapolation]%
\label{thm:geom_extrapolation}
	\lean{MPSTensor.SectorWeightData.power_sums_eq_of_eventually_eq}
	\leanok
	If two finite families of nonzero complex numbers have equal power sums
	$\sum_i \alpha_i^k = \sum_j \beta_j^k$ for all $k > N_0$, then
	equality holds for all $k \geq 0$.
	The proof uses induction on the number of terms with a telescoping
	argument: subtract a geometric factor to reduce the number of terms,
	apply the induction hypothesis, and conclude via a geometric
	recurrence.
\end{theorem}

\begin{theorem}[Route~B equal-case corollary: weight multiset recovery]%
\label{thm:route_b_equal}
	\lean{MPSTensor.SectorWeightData.weight_multiset_eq_of_sameMPV_bnt}
	\leanok
	\uses{thm:coeff_eventually_eq, thm:geom_extrapolation, thm:power_sum_multiset, def:paper_sector_data, def:bnt}
	Let $S$ and $T$ be two paper sector data over the same basis with
	the same multiplicities ($r_j^S = r_j^T$ for all~$j$).
	If the basis is eventually linearly independent and the total MPVs
	of $S$ and $T$ agree, then for each basis block~$j$ the sector weight
	multisets $\{\mu_{j,q}^S\}$ and $\{\mu_{j,q}^T\}$ are equal.

	The proof chains:
	\begin{enumerate}
		\item BNT LI $+$ equal MPVs $\Rightarrow$ eventual coefficient
		      equality (Theorem~\ref{thm:coeff_eventually_eq});
		\item Geometric extrapolation $\Rightarrow$ full coefficient
		      equality for all positive~$k$
		      (Theorem~\ref{thm:geom_extrapolation});
		\item Newton--Girard $\Rightarrow$ equal weight multisets
		      (Theorem~\ref{thm:power_sum_multiset}).
	\end{enumerate}
\end{theorem}

\begin{remark}[Current formalized endpoint versus the literature]
	The formalization provides three complementary equal-case results:
	\begin{enumerate}
		\item Theorem~\ref{thm:ft_equal_explicit}: the explicit-coefficient
		      equal-case theorem (upgrade of
		      Theorem~\ref{thm:ft_proportional} with $c_N = 1$),
		      concluding a global gauge equivalence of the reindexed
		      weighted block-diagonal tensors.
		\item Theorem~\ref{thm:ft_equal_same_structure}: the common
		      block-structure special case.
		\item Theorem~\ref{thm:route_b_equal}: the Route~B equal-case
		      corollary, which from the paper's richer multiplicity data
		      recovers the sector weight multisets without requiring
		      explicit coefficient convergence hypotheses.
	\end{enumerate}
	The literature-level unconditional Theorem~\ref{thm:ft_equal}
	(\cite[Corollary~IV.5]{Cirac2021Matrix}) remains \texttt{\textbackslash notready}.
	To close it, one must compose the Route~B weight recovery with:
	(a)~the proportional FT for block matching and gauge-phase equivalences,
	(b)~phase absorption into weights, and
	(c)~global gauge construction via permutation and cast transport
	(all of which are now formalized individually in \texttt{Multi.lean}).
\end{remark}

\section{Periodic Fundamental Theorem (de las Cuevas--Cirac--Schuch--Pérez-García)}%
\label{sec:periodic_ft}

The theorems in the preceding sections all require that each block has a
\emph{primitive} transfer map, i.e.\ peripheral spectrum exactly $\{1\}$.
This is achieved by blocking the original tensor by a common period.
The paper \cite{DeLasCuevas2017Irreducible} develops a framework that
avoids this blocking step, giving a fundamental theorem that applies
directly to periodic tensors.

\subsection{Irreducible form}

\begin{definition}[Periodic tensor]
	\label{def:periodic_tensor}
	\lean{MPSTensor.IsPeriodic}
	\leanok
	An MPS tensor $A$ of bond dimension $D$ is
	\emph{$m$-periodic} (for some $m \ge 1$) if the following conditions hold:
	\begin{enumerate}
		\item $A$ is irreducible (no nontrivial invariant projection),
		\item $A$ is left-canonical (trace-preserving normalization:
		      $\sum_i (A^i)^\dagger A^i = \Id_D$),
		\item $m \ge 1$ (the period is positive),
		\item the peripheral spectrum of $\E_A$ is exactly the set of
		      $m$-th roots of unity
		      $\{e^{2\pi i k/m} : 0 \le k < m\}$, and
		\item there exists a primitive $m$-th root of unity.
	\end{enumerate}
	A $1$-periodic tensor is exactly an irreducible, left-canonical tensor
	with primitive transfer map (Definition~\ref{def:primitive_channel}).
\end{definition}

\begin{definition}[Irreducible form]
	\label{def:irreducible_form}
	\lean{MPSTensor.IsIrreducibleForm}
	\leanok
	A tensor $A$ is in \emph{irreducible form} if it is block-diagonal
	\[
		A^i = \bigoplus_{j \in J} \mu_j\, A_j^i,
	\]
	where each block $A_j$ is an $m_j$-periodic tensor,
	each weight $\mu_j$ has strictly positive real part
	($\Re(\mu_j) > 0$), and the reassembled block-diagonal tensor
	generates the same MPV family as $A$.
	The pairwise non-repetition condition (no two distinct blocks are
	repeated versions of each other) is bundled separately in the
	\texttt{BasisOfPeriodicTensors} structure rather than in
	\texttt{IsIrreducibleForm}.
	This is the standard form introduced in \cite[\S II]{DeLasCuevas2017Irreducible}.
	It extends the normal canonical form
	(Definition~\ref{def:is_normal_canonical_form}),
	which requires the stronger condition that every block is $1$-periodic.
\end{definition}

\begin{remark}[Relation to normal canonical form]\notready
	\label{rem:irr_form_vs_ncf}
	Every tensor in normal canonical form
	(Definition~\ref{def:is_normal_canonical_form})
	is in irreducible form with all periods equal to~$1$.
	The irreducible form is strictly more general: it accommodates
	blocks such as the antiferromagnet, which are $2$-periodic.

	The current formalization (Chapters~\ref{ch:canonical}--\ref{ch:assembly})
	works exclusively with period-$1$ blocks, obtaining them by a blocking step.
	Extending the formalization to the irreducible form of
	\cite{DeLasCuevas2017Irreducible} would require:
	\begin{enumerate}
		\item a definition of $m$-periodic irreducible tensors and
		      their cyclic-sector decomposition (partial infrastructure
		      exists in Section~\ref{sec:cyclic_sectors} of
		      Chapter~\ref{ch:canonical}),
		\item the off-diagonal structure of a periodic block's transfer
		      matrix (the cyclic permutation of sectors), and
		\item the notion of equivalence between periodic blocks that
		      accounts for the $Z$-gauge freedom described below.
	\end{enumerate}
\end{remark}

\subsection{The \texorpdfstring{$Z$}{Z}-gauge degree of freedom}

\begin{definition}[Entrywise periodic \(Z\)-gauge ratio]\label{def:z_gauge_entry}
	\lean{MPSTensor.zGaugeEntry}
	\leanok
	Given complex weights $\mu,\nu$, define the entrywise
	\(Z\)-gauge ratio by
	\[
		z(\mu,\nu) := \mu/\nu.
	\]
\end{definition}

\begin{lemma}[Matched powers imply root-of-unity ratio]\label{lem:z_gauge_entry_pow}
	\lean{MPSTensor.zGaugeEntry_pow_eq_one_of_pow_eq}
	\leanok
	\uses{def:z_gauge_entry}
	If $\mu^m=\nu^m$ and $\nu\neq 0$, then
	\[
		z(\mu,\nu)^m = 1.
	\]
\end{lemma}

\begin{lemma}[Ratio rescales denominator back to numerator]\label{lem:z_gauge_entry_mul}
	\lean{MPSTensor.zGaugeEntry_mul_right}
	\leanok
	\uses{def:z_gauge_entry}
	If $\nu\neq 0$, then
	\[
		z(\mu,\nu)\,\nu=\mu.
	\]
\end{lemma}

\begin{definition}[Diagonal periodic \(Z\)-gauge matrix]\label{def:z_gauge_diagonal}
	\lean{MPSTensor.zGaugeDiagonal}
	\leanok
	\uses{def:z_gauge_entry}
	For matched weight families $(\mu_i)_i$ and $(\nu_i)_i$ on a finite
	index type, define the diagonal matrix
	\[
		Z := \diag\!\bigl(z(\mu_i,\nu_i)\bigr)_i.
	\]
\end{definition}

\begin{lemma}[Diagonal \(Z\)-gauge has finite order under matched powers]\label{lem:z_gauge_diagonal_pow}
	\lean{MPSTensor.zGaugeDiagonal_pow_eq_one}
	\leanok
	\uses{def:z_gauge_diagonal, lem:z_gauge_entry_pow}
	If $\mu_i^m=\nu_i^m$ and $\nu_i\neq 0$ for all~$i$, then
	\[
		Z^m=\Id.
	\]
\end{lemma}

\begin{lemma}[Diagonal \(Z\)-gauge transports diagonal weights]\label{lem:z_gauge_diagonal_mul}
	\lean{MPSTensor.zGaugeDiagonal_mul_diagonal}
	\leanok
	\uses{def:z_gauge_diagonal, lem:z_gauge_entry_mul}
	Under $\nu_i\neq 0$ for all~$i$,
	\[
		Z\,\diag(\nu)=\diag(\mu).
	\]
\end{lemma}

\begin{definition}[$Z$-gauge equivalence]
	\label{def:z_gauge_equiv}
	\lean{MPSTensor.ZGaugeEquiv}
	\leanok
	Let $A$ be an $m$-periodic irreducible tensor of bond dimension $D$.
	A \emph{$Z$-gauge transformation} of $A$ is a diagonal unitary
	matrix $Z \in \MN{D}$ satisfying:
	\begin{enumerate}
		\item $Z$ commutes with every Kraus operator: $[A^i, Z] = 0$
		      for all physical indices $i$, and
		\item $Z^m = \Id_D$ (all diagonal entries are $m$-th roots of unity).
	\end{enumerate}
	Replacing $A$ by $ZA$ (i.e.\ $A^i \mapsto Z A^i$) leaves the
	MPV family invariant for lengths $N$ that are multiples of $m$,
	because the factor $\tr(Z^N) = \tr(\Id) = D$ for $N \equiv 0 \pmod{m}$.
	Two tensors $A, B$ are \emph{$Z$-gauge equivalent} if there exists
	an invertible $Y$ and a $Z$-gauge transformation $Z$ for $A$ such that
	$Z A^i = Y B^i Y^{-1}$ for all~$i$.
\end{definition}

\begin{remark}[$Z$-gauge is trivial for period-$1$ blocks]
	\label{rem:z_gauge_trivial_period1}
	\lean{MPSTensor.ZGaugeEquiv.of_period_one}
	\leanok
	When $m = 1$ the condition $Z^1 = \Id$ forces $Z = \Id$,
	so $Z$-gauge equivalence reduces to ordinary gauge equivalence
	(Definition~\ref{def:gauge_equiv}).
	The $Z$-gauge is therefore a genuine new degree of freedom
	that appears only for $m > 1$.
\end{remark}

\subsection{Statement of the periodic Fundamental Theorem}

The following is Theorem~3.8 (``equal case'') of \cite{DeLasCuevas2017Irreducible}.

\begin{theorem}[Periodic Fundamental Theorem]\notready
	\label{thm:periodic_ft}
	\uses{def:irreducible_form, def:z_gauge_equiv}
	Let $A$ and $B$ be MPS tensors in irreducible form,
	with bases of periodic blocks
	$\{A_j\}_{j \in J}$ and $\{B_k\}_{k \in K}$ respectively.
	Suppose that $\mathcal{V}(A) = \mathcal{V}(B)$
	(equal MPV families for all lengths $N$).

	Then $|J| = |K|$ and there is a bijection $\pi : J \to K$
	such that $A_j$ and $B_{\pi(j)}$ are $m_j$-periodic for the same
	period $m_j$.
	Moreover, for each $j \in J$ there exists:
	\begin{enumerate}
		\item an invertible matrix $Y_j$ of size $D_j \times D_{\pi(j)}$,
		      and
		\item a diagonal unitary matrix $Z_j$ of size $D_j \times D_j$
		      satisfying $Z_j^{m_j} = \Id$ and
		      $[A_j^i, Z_j] = 0$ for all~$i$,
	\end{enumerate}
	such that
	\[
		Z_j\, A_j^i = Y_j\, B_{\pi(j)}^i\, Y_j^{-1}
		\quad\text{for all physical indices } i.
	\]
	Equivalently, setting $Y = \bigoplus_j Y_j \otimes \Id_{r_j}$
	and $Z = \bigoplus_j Z_j \otimes \Id_{r_j}$, one has
	$Z A^i = Y B^i Y^{-1}$ for all $i$,
	with $[A^i, Z] = 0$ and $\mathcal{V}(A) = \mathcal{V}(ZA)$.
\end{theorem}

\begin{remark}[Comparison with the 1606.00608 approach]\notready
	\label{rem:periodic_vs_blocking_ft}
	Theorem~\ref{thm:periodic_ft} and the equal-MPV Fundamental Theorem
	(Theorem~\ref{thm:ft_equal}) both characterize the gauge freedom
	between two tensors with equal MPV families, but they do so at
	different levels:

	\begin{enumerate}
		\item \textbf{Blocking approach} (Theorem~\ref{thm:ft_equal},
		      following \cite{Cirac2017MPDO_AnnPhys}):
		      Block both $A$ and $B$ by a common period $p$ so that all
		      blocks of $A^{[p]}$ and $B^{[p]}$ are primitive
		      ($1$-periodic). Apply the aperiodic fundamental theorem to
		      the blocked tensors. The conclusion is that $A^{[p]}$ and
		      $B^{[p]}$ are gauge equivalent. The $Z$-gauge freedom
		      does not appear because blocking absorbs the periodicity.

		\item \textbf{Periodic approach} (Theorem~\ref{thm:periodic_ft},
		      following \cite{DeLasCuevas2017Irreducible}):
		      Work directly with the periodic blocks of $A$ and $B$.
		      The conclusion includes an extra $Z$-gauge transformation
		      reflecting the non-trivial phase ambiguity that arises for
		      periodic blocks. The $Z$ matrix satisfies $Z^m = \Id$ and
		      commutes with all Kraus operators; it is invisible to the
		      MPV family because it only introduces $m$-th roots of unity
		      in the trace.
	\end{enumerate}

	In the current formalization, only the blocking approach has been
	implemented. Formalizing the periodic approach would require
	developing the irreducible form and $Z$-gauge infrastructure
	outlined in Remark~\ref{rem:irr_form_vs_ncf}.
\end{remark}

\begin{remark}[Proportional periodic FT]\notready
	\label{rem:periodic_ft_proportional}
	There is also a proportional-case version of the periodic Fundamental Theorem
	(Theorem~3.4, ``proportional case'', in \cite{DeLasCuevas2017Irreducible}):
	if $\mathcal{V}(A) = \lambda_N \mathcal{V}(B)$ for scalars $\lambda_N$ with
	$\lambda_N \to \lambda \neq 0$, then the bases of periodic tensors of $A$
	and $B$ are equivalent (up to the same gauge and $Z$ freedom as above).
	This extends Theorems~\ref{thm:ft_proportional}
	and~\ref{thm:ft_proportional_nt} to the periodic setting.
\end{remark}

\subsection{Structural FT after blocking}

\begin{theorem}[Structural Fundamental Theorem after blocking]%
	\label{thm:ft_after_blocking_structural}
	\lean{MPSTensor.fundamentalTheorem_after_blocking_1606_structural}
	\leanok
	\uses{def:periodic_tensor}
	For any two MPS tensors $A$ and $B$ generating the same MPV family
	($\mathcal{V}(A) = \mathcal{V}(B)$), there exist blocking periods
	$p_A, p_B \ge 1$ such that both blocked tensors $A^{[p_A]}$ and
	$B^{[p_B]}$ admit TP-primitive block-diagonal decompositions:
	\[
		A^{[p_A]} \simeq \bigoplus_{j=1}^{r_A} \mu_j^A\, A_j,
		\qquad
		B^{[p_B]} \simeq \bigoplus_{k=1}^{r_B} \mu_k^B\, B_k,
	\]
	where each $A_j$, $B_k$ is left-canonical with primitive transfer map,
	every weight is nonzero, and every block has positive bond dimension.
\end{theorem}

\begin{proof}\leanok
	Apply the blocking-based reduction
	(\texttt{exists\_tp\_primitive\_blockDecomp\_after\_blocking})
	independently to $A$ and~$B$.
\end{proof}

\begin{remark}[Remaining steps for the full 1606.00608 FT]
	\label{rem:ft_1606_remaining}
	Theorem~\ref{thm:ft_after_blocking_structural} is the current formalized
	endpoint toward the full end-to-end Fundamental Theorem of
	\cite{Cirac2017MPDO_AnnPhys}. The remaining steps are:
	\begin{enumerate}
		\item \textbf{Common blocking period}: Re-block both decompositions
		      to $p = \mathrm{lcm}(p_A, p_B)$, requiring iterated-blocking
		      infrastructure.
		\item \textbf{Sector irreducibility}: Each cyclic sector of a blocked
		      periodic block should be irreducible (see~\#242).
		\item \textbf{Normal form per sector}: Already proved
		      (\texttt{isNormal\_of\_tp\_primitive\_irreducible}).
		\item \textbf{BNT separation + weight ordering + gauge-phase matching}:
		      Already proved individually; needs combination after steps~1--2
		      (see~\#243).
	\end{enumerate}
\end{remark}
